\title{2022 年高考数学 浙江卷}
\date{}
\maketitle

\section*{一、选择题}
本大题共10小题，每小题4分，共40分．在每小题给出的四个选项中，只有一项是符合题目要求的．

\begin{question}[points=4]
设集合 $A = \{1,2\}, B = \{2,4,6\}$ ，则 $A \cup B =$ （ ）
\begin{choices}
  \item $\{2\}$
  \item $\{1,2\}$
  \item $\{2,4,6\}$
  \item $\{1,2,4,6\}$
\end{choices}
\begin{solution}
\textbf{答案：}D

\textbf{分析：}利用并集的定义可得正确的选项.

\textbf{详解：}$A \cup B = \{1,2,4,6\}$，故选D.
\end{solution}
\end{question}

\begin{question}[points=4]
已知 $a,b \in \mathbb{R}, a+3i = (b+i)i$ （$i$ 为虚数单位），则（ ）
\begin{choices}
  \item $a = 1, b = -3$
  \item $a = -1, b = 3$
  \item $a = -1, b = -3$
  \item $a = 1, b = 3$
\end{choices}
\begin{solution}
\textbf{答案：}B

\textbf{分析：}利用复数相等的条件可求 $a,b$.

\textbf{详解：}由 $a+3i = (b+i)i = -1 + bi$，得 $a = -1, b = 3$，故选B.
\end{solution}
\end{question}

\begin{question}[points=4]
若实数 $x,y$ 满足约束条件 $\begin{cases} x-2 \ge 0, \\ 2x+y-7 \le 0, \\ x-y-2 \le 0, \end{cases}$ 则 $z = 3x+4y$ 的最大值是（ ）
\begin{choices}
  \item $20$
  \item $18$
  \item $13$
  \item $6$
\end{choices}
\begin{solution}
\textbf{答案：}B

\textbf{分析：}在平面直角坐标系中画出可行域，平移动直线 $z=3x+4y$ 后可求最大值.

\textbf{详解：}不等式组对应的可行域如图所示：当动直线 $3x+4y-z=0$ 过 $A$ 时 $z$ 有最大值．由 $\begin{cases} x=2 \\ 2x+y-7=0 \end{cases}$ 得 $A(2,3)$，故 $z_{\max}=3\times2+4\times3=18$，故选B.
\end{solution}
\end{question}

\begin{question}[points=4]
设 $x \in \mathbb{R}$，则 $\sin x = 1$ 是 $\cos x = 0$ 的（ ）
\begin{choices}
  \item 充分不必要条件
  \item 必要不充分条件
  \item 充分必要条件
  \item 既不充分也不必要条件
\end{choices}
\begin{solution}
\textbf{答案：}A

\textbf{分析：}由三角函数的性质结合充分条件、必要条件的定义即可得解.

\textbf{详解：}因为 $\sin^2 x + \cos^2 x = 1$，当 $\sin x = 1$ 时，$\cos x = 0$，充分性成立；当 $\cos x = 0$ 时，$\sin x = \pm 1$，必要性不成立；所以 $\sin x = 1$ 是 $\cos x = 0$ 的充分不必要条件，故选A.
\end{solution}
\end{question}

\begin{question}[points=4]
某几何体的三视图如图所示（单位：$\mathrm{cm}$），则该几何体的体积（单位：$\mathrm{cm}^3$）是（ ）

\begin{center}
  \includegraphics[width=.38\linewidth]{question-5-1.jpg}
\end{center}
\begin{choices}
  \item $22\pi$
  \item $8\pi$
  \item $\dfrac{22}{3}\pi$
  \item $\dfrac{16}{3}\pi$
\end{choices}
\begin{solution}
\textbf{答案：}C

\textbf{分析：}根据三视图还原几何体可知，原几何体是一个半球，一个圆柱，一个圆台组合成的几何体，即可根据球，圆柱，圆台的体积公式求出．

\textbf{详解：}由三视图可知，该几何体是一个半球，一个圆柱，一个圆台组合成的几何体，球的半径、圆柱的底面半径、圆台的上底面半径都为 $1\mathrm{cm}$，圆台的下底面半径为 $2\mathrm{cm}$，所以体积 $V = \frac{1}{2} \times \frac{4}{3}\pi \times 1^3 + \pi \times 1^2 \times 2 + \frac{1}{3} \times 2 \times [\pi \times 2^2 + \pi \times 1^2 + \sqrt{\pi \times 2^2 \times \pi \times 1^2}] = \frac{22}{3}\pi \mathrm{cm}^3$，故选C.
\end{solution}
\end{question}

\begin{question}[points=4]
为了得到函数 $y = 2\sin 3x$ 的图象，只要把函数 $y = 2\sin\left(3x + \frac{\pi}{5}\right)$ 图象上所有的点（ ）
\begin{choices}
  \item 向左平移 $\dfrac{\pi}{5}$ 个单位长度
  \item 向右平移 $\dfrac{\pi}{5}$ 个单位长度
  \item 向左平移 $\dfrac{\pi}{15}$ 个单位长度
  \item 向右平移 $\dfrac{\pi}{15}$ 个单位长度
\end{choices}
\begin{solution}
\textbf{答案：}D

\textbf{分析：}根据三角函数图象的变换法则即可求出．

\textbf{详解：}$y = 2\sin 3x = 2\sin\left[3\left(x - \frac{\pi}{15}\right) + \frac{\pi}{5}\right]$，所以向右平移 $\frac{\pi}{15}$ 个单位长度，故选D.
\end{solution}
\end{question}

\begin{question}[points=4]
已知 $2^a = 5, \log_8 3 = b$，则 $4^{a-3b} =$ （ ）
\begin{choices}
  \item $25$
  \item $5$
  \item $\dfrac{25}{9}$
  \item $\dfrac{5}{3}$
\end{choices}
\begin{solution}
\textbf{答案：}C

\textbf{分析：}根据指数式与对数式的互化，幂的运算性质以及对数的运算性质即可解出．

\textbf{详解：}因为 $2^a = 5$，$b = \log_8 3 = \frac{1}{3} \log_2 3$，所以 $2^{3b} = 3$，则 $4^{a-3b} = \frac{4^a}{4^{3b}} = \frac{(2^a)^2}{(2^{3b})^2} = \frac{5^2}{3^2} = \frac{25}{9}$，故选C.
\end{solution}
\end{question}

\begin{question}[points=4]
如图，已知正三棱柱 $ABC-A_1B_1C_1, AC = AA_1$，$E,F$ 分别是棱 $BC, A_1C_1$ 上的点．记 $EF$ 与 $AA_1$ 所成的角为 $\alpha$，$EF$ 与平面 $ABC$ 所成的角为 $\beta$，二面角 $F-BC-A$ 的平面角为 $\gamma$，则（ ）
\begin{choices}
  \item $\alpha \le \beta \le \gamma$
  \item $\beta \le \alpha \le \gamma$
  \item $\beta \le \gamma \le \alpha$
  \item $\alpha \le \gamma \le \beta$
\end{choices}
\begin{solution}
\textbf{答案：}A

\textbf{分析：}先用几何法表示出 $\alpha,\beta,\gamma$，再根据边长关系即可比较大小．

\textbf{详解：}如图所示，过点 $F$ 作 $FP \perp AC$ 于 $P$，过 $P$ 作 $PM \perp BC$ 于 $M$，连接 $PE$，则 $\alpha = \angle EFP$，$\beta = \angle FEP$，$\gamma = \angle FMP$，$\tan\alpha = \frac{PE}{FP} \le 1$，$\tan\beta = \frac{FP}{PE} \ge 1$，$\tan\gamma = \frac{FP}{PM} \ge \frac{FP}{PE} = \tan\beta$，所以 $\alpha \le \beta \le \gamma$，故选A.
\end{solution}
\end{question}

\begin{question}[points=4]
已知 $a,b \in \mathbb{R}$，若对任意 $x \in \mathbb{R}, a|x-b| + |x-4| - |2x-5| \ge 0$，则（ ）
\begin{choices}
  \item $a \le 1, b \ge 3$
  \item $a \le 1, b \le 3$
  \item $a \ge 1, b \ge 3$
  \item $a \ge 1, b \le 3$
\end{choices}
\begin{solution}
\textbf{答案：}D

\textbf{分析：}将问题转换为 $a|x-b| \ge |2x-5| - |x-4|$，再结合画图求解．

\textbf{详解：}设 $f(x)=a|x-b|$，$g(x)=|2x-5|-|x-4| = \begin{cases} 1-x, & x \le \frac{5}{2} \\ 3x-9, & \frac{5}{2} < x < 4 \\ x-1, & x \ge 4 \end{cases}$，则 $f(x)$ 的图象恒在 $g(x)$ 的上方（可重合），由图可知，$a \ge 3, 1 \le b \le 3$，或 $1 \le a < 3, 1 \le b \le 4-\frac{3}{a} \le 3$，故选D.
\end{solution}
\end{question}

\begin{question}[points=4]
已知数列 $\{a_n\}$ 满足 $a_1 = 1, a_{n+1} = a_n - \frac{1}{3}a_n^2$ $(n \in \mathbb{N}^*)$，则（ ）
\begin{choices}
  \item $2 < 100a_{100} < \dfrac{5}{2}$
  \item $\dfrac{5}{2} < 100a_{100} < 3$
  \item $3 < 100a_{100} < \dfrac{7}{2}$
  \item $\dfrac{7}{2} < 100a_{100} < 4$
\end{choices}
\begin{solution}
\textbf{答案：}B

\textbf{分析：}先通过递推关系式确定 $\{a_n\}$ 除去 $a_1$，其他项都在 $(0,1)$ 范围内，再利用递推公式变形得到 $\frac{1}{a_{n+1}} - \frac{1}{a_n} = \frac{1}{3-a_n} > \frac{1}{3}$，累加可求出 $\frac{1}{a_n} > \frac{1}{3}(n+2)$，得出 $100a_{100} < 3$；再利用 $\frac{1}{a_{n+1}} - \frac{1}{a_n} = \frac{1}{3-a_n} < \frac{1}{3-\frac{3}{n+2}} = \frac{1}{3}\left(1+\frac{1}{n+1}\right)$，累加可求出 $\frac{1}{a_n} - 1 < \frac{1}{3}(n-1) + \frac{1}{3}\left(\frac{1}{2}+\frac{1}{3}+\cdots+\frac{1}{n}\right)$，再次放缩可得出 $100a_{100} > \frac{5}{2}$．

\textbf{详解：}由 $a_1=1$ 易得 $a_2=\frac{2}{3} \in (0,1)$，依次类推 $a_n \in (0,1)$．由 $a_{n+1}=a_n\left(1-\frac{1}{3}a_n\right)$ 得 $\frac{1}{a_{n+1}} = \frac{3}{a_n(3-a_n)} = \frac{1}{a_n} + \frac{1}{3-a_n}$，所以 $\frac{1}{a_{n+1}} - \frac{1}{a_n} = \frac{1}{3-a_n} > \frac{1}{3}$，累加得 $\frac{1}{a_n} - 1 > \frac{1}{3}(n-1)$，即 $\frac{1}{a_n} > \frac{1}{3}(n+2)$，所以 $a_n < \frac{3}{n+2}$，$100a_{100} < \frac{100}{34} < 3$．又 $\frac{1}{a_{n+1}} - \frac{1}{a_n} = \frac{1}{3-a_n} < \frac{1}{3-\frac{3}{n+2}} = \frac{1}{3}\left(1+\frac{1}{n+1}\right)$，累加得 $\frac{1}{a_{100}} - 1 < \frac{1}{3}(99) + \frac{1}{3}\left(\frac{1}{2}+\cdots+\frac{1}{99}\right) < 33 + \frac{1}{3}\left(\frac{1}{2}\times4 + \frac{1}{6}\times94\right) < 39$，即 $\frac{1}{a_{100}} < 40$，$a_{100} > \frac{1}{40}$，所以 $100a_{100} > \frac{5}{2}$，故选B.
\end{solution}
\end{question}

\section*{二、填空题}
本大题共7小题，单空题每题4分，多空题每空3分，共36分．

\begin{question}[points=4]
我国南宋著名数学家秦九韶，发现了从三角形三边求面积的公式，他把这种方法称为“三斜求积”，它填补了我国传统数学的一个空白．如果把这个方法写成公式，就是 $S = \frac{1}{4}\left[c^2 a^2 - \left(\frac{c^2 + a^2 - b^2}{2}\right)^2\right]$，其中 $a,b,c$ 是三角形的三边，$S$ 是三角形的面积．设某三角形的三边 $a = \sqrt{2}, b = 3, c = 2$，则该三角形的面积 $S =$ ．
\fillin[{$\dfrac{\sqrt{23}}{4}$}]
\begin{solution}
\textbf{答案：}$\dfrac{\sqrt{23}}{4}$

\textbf{分析：}根据题中所给的公式代值解出．

\textbf{详解：}$S = \frac{1}{4}\left[4\times2 - \left(\frac{4+2-3}{2}\right)^2\right] = \frac{\sqrt{23}}{4}$
\end{solution}
\end{question}

\begin{question}[points=6]
已知多项式 $(x+2)(x-1)^4 = a_0 + a_1 x + a_2 x^2 + a_3 x^3 + a_4 x^4 + a_5 x^5$，则 $a_2 =$ ，$a_1 + a_2 + a_3 + a_4 + a_5 =$ ．
\fillin[{$8$；$-2$}]
\begin{solution}
\textbf{答案：}$8$；$-2$

\textbf{分析：}第一空利用二项式定理直接求解即可，第二空赋值去求，令 $x=0$ 求出 $a_0$，再令 $x=1$ 即可得出答案．

\textbf{详解：}含 $x^2$ 的项为 $x\cdot C_4^3 x \cdot (-1)^3 + 2\cdot C_4^2 x^2 \cdot (-1)^2 = -4x^2+12x^2=8x^2$，故 $a_2=8$．令 $x=0$ 得 $2=a_0$，令 $x=1$ 得 $0=a_0+a_1+\cdots+a_5$，所以 $a_1+\cdots+a_5 = -2$.
\end{solution}
\end{question}

\begin{question}[points=6]
若 $3\sin\alpha - \sin\beta = \sqrt{10}$，$\alpha+\beta = \dfrac{\pi}{2}$，则 $\sin\alpha =$ ，$\cos 2\beta =$ ．
\fillin[{$\dfrac{3\sqrt{10}}{10}$；$\dfrac{4}{5}$}]
\begin{solution}
\textbf{答案：}$\dfrac{3\sqrt{10}}{10}$；$\dfrac{4}{5}$

\textbf{分析：}先通过诱导公式变形，得到 $\alpha$ 的同角等式关系，再利用辅助角公式化简成正弦型函数方程，可求出 $\alpha$，接下来再求 $\beta$．

\textbf{详解：}由 $\alpha+\beta=\frac{\pi}{2}$ 得 $\sin\beta = \cos\alpha$，代入得 $3\sin\alpha - \cos\alpha = \sqrt{10}$，即 $\sqrt{10}\sin(\alpha-\theta)=\sqrt{10}$，其中 $\sin\theta=\frac{1}{\sqrt{10}}$，$\cos\theta=\frac{3}{\sqrt{10}}$，所以 $\alpha-\theta = \frac{\pi}{2}+2k\pi$，$\sin\alpha = \sin(\theta+\frac{\pi}{2}) = \cos\theta = \frac{3\sqrt{10}}{10}$，$\cos2\beta = 2\cos^2\beta-1 = 2\sin^2\alpha-1 = \frac{4}{5}$.
\end{solution}
\end{question}

\begin{question}[points=6]
已知函数 $f(x) = \begin{cases} -x^2+2, & x \le 1, \\ x+\dfrac{1}{x}-1, & x > 1, \end{cases}$ 则 $f\left(f\left(\frac{1}{2}\right)\right) =$ ；若当 $x \in [a,b]$ 时，$1 \le f(x) \le 3$，则 $b-a$ 的最大值是．
\fillin[{$\dfrac{37}{28}$；$3+\sqrt{3}$}]
\begin{solution}
\textbf{答案：}$\dfrac{37}{28}$；$3+\sqrt{3}$

\textbf{分析：}结合分段函数的解析式求函数值，由条件求出 $a$ 的最小值，$b$ 的最大值即可.

\textbf{详解：}$f(\frac{1}{2}) = -(\frac{1}{2})^2+2 = \frac{7}{4}$，$f(\frac{7}{4}) = \frac{7}{4}+\frac{4}{7}-1 = \frac{37}{28}$．当 $x\le1$ 时，由 $1\le -x^2+2 \le 3$ 得 $-1\le x\le1$；当 $x>1$ 时，由 $1\le x+\frac{1}{x}-1\le 3$ 得 $1<x\le 2+\sqrt{3}$；所以 $[a,b]\subseteq[-1,2+\sqrt{3}]$，$b-a$ 最大值为 $3+\sqrt{3}$.
\end{solution}
\end{question}

\begin{question}[points=6]
现有7张卡片，分别写上数字1,2,2,3,4,5,6．从这7张卡片中随机抽取3张，记所抽取卡片上数字的最小值为 $\xi$，则 $P(\xi=2) =$ ，$E(\xi) =$ ．
\fillin[{$\dfrac{16}{35}$；$\dfrac{12}{7}$}]
\begin{solution}
\textbf{答案：}$\dfrac{16}{35}$；$\dfrac{12}{7}$

\textbf{分析：}利用古典概型概率公式求 $P(\xi=2)$，由条件求 $\xi$ 分布列，再由期望公式求其期望.

\textbf{详解：}从7张中任取3张共有 $C_7^3=35$ 种取法，$\xi=2$ 时，抽取的卡片中不含1且至少有一张2，取法有 $C_4^1 + C_2^1 C_4^2 = 4+12=16$，所以 $P(\xi=2)=\frac{16}{35}$；$\xi$ 的取值为1,2,3,4，$P(\xi=1)=\frac{C_6^2}{C_7^3}=\frac{15}{35}$，$P(\xi=2)=\frac{16}{35}$，$P(\xi=3)=\frac{C_3^2}{C_7^3}=\frac{3}{35}$，$P(\xi=4)=\frac{1}{35}$，所以 $E(\xi)=1\times\frac{15}{35}+2\times\frac{16}{35}+3\times\frac{3}{35}+4\times\frac{1}{35}=\frac{12}{7}$.
\end{solution}
\end{question}

\begin{question}[points=4]
已知双曲线 $\dfrac{x^2}{a^2} - \dfrac{y^2}{b^2} = 1$ $(a>0,b>0)$ 的左焦点为 $F$，过 $F$ 且斜率为 $\dfrac{b}{4a}$ 的直线交双曲线于点 $A(x_1,y_1)$，交双曲线的渐近线于点 $B(x_2,y_2)$ 且 $x_1<0<x_2$．若 $|FB| = 3|FA|$，则双曲线的离心率是．
\fillin[{$\dfrac{3\sqrt{6}}{4}$}]
\begin{solution}
\textbf{答案：}$\dfrac{3\sqrt{6}}{4}$

\textbf{分析：}联立直线 $AB$ 和渐近线 $l_2: y=\frac{b}{a}x$ 方程，可求出点 $B$，再根据 $|FB|=3|FA|$ 可求得点 $A$，最后根据点 $A$ 在双曲线上，即可解出离心率．

\textbf{详解：}直线 $AB: y = \frac{b}{4a}(x+c)$，渐近线 $l_2: y = \frac{b}{a}x$，联立得 $B\left(\frac{c}{3},\frac{bc}{3a}\right)$，由 $|FB|=3|FA|$ 得 $A\left(-\frac{5c}{9},\frac{bc}{9a}\right)$，代入双曲线方程得 $\frac{25c^2}{81a^2} - \frac{b^2c^2}{81a^2 b^2} = 1$，解得 $\frac{c^2}{a^2} = \frac{81}{24}$，所以 $e = \frac{3\sqrt{6}}{4}$.
\end{solution}
\end{question}

\begin{question}[points=4]
设点 $P$ 在单位圆的内接正八边形 $A_1 A_2 \cdots A_8$ 的边 $A_1 A_2$ 上，则 $\overrightarrow{PA_1}^2 + \overrightarrow{PA_2}^2 + \cdots + \overrightarrow{PA_8}^2$ 的取值范围是．

\begin{center}
  \includegraphics[width=.38\linewidth]{question-17-1.jpg}
\end{center}
\fillin[{$[12+2\sqrt{2}, 16]$}]
\begin{solution}
\textbf{答案：}$[12+2\sqrt{2}, 16]$

\textbf{分析：}根据正八边形的结构特征，分别以圆心为原点，$A_7A_3$ 所在直线为 $x$ 轴，$A_5A_1$ 所在直线为 $y$ 轴建立平面直角坐标系，即可求出各顶点的坐标，设 $P(x,y)$，再根据平面向量模的坐标计算公式即可得到 $\overrightarrow{PA_1}^2+\cdots+\overrightarrow{PA_8}^2 = 8(x^2+y^2)+8$，然后利用 $\cos 22.5^\circ \le |OP| \le 1$ 即可解出．

\textbf{详解：}建系如图，各顶点坐标：$A_1(0,1), A_2(\frac{\sqrt{2}}{2},\frac{\sqrt{2}}{2}), A_3(1,0), A_4(\frac{\sqrt{2}}{2},-\frac{\sqrt{2}}{2}), A_5(0,-1), A_6(-\frac{\sqrt{2}}{2},-\frac{\sqrt{2}}{2}), A_7(-1,0), A_8(-\frac{\sqrt{2}}{2},\frac{\sqrt{2}}{2})$．设 $P(x,y)$，则 $\sum PA_i^2 = 8(x^2+y^2)+8$．$|OP|$ 的最小值为 $\cos22.5^\circ = \frac{\sqrt{2+\sqrt{2}}}{2}$，最大值为1，所以 $x^2+y^2 \in [\frac{1+\cos45^\circ}{2},1] = [\frac{2+\sqrt{2}}{4},1]$，故取值范围为 $[12+2\sqrt{2},16]$.
\end{solution}
\end{question}

\section*{三、解答题}
本大题共5小题，共74分．解答应写出文字说明、证明过程或演算步骤．

\begin{question}[points=14]
在 $\triangle ABC$ 中，角 $A,B,C$ 所对的边分别为 $a,b,c$．已知 $4a = \sqrt{5}c, \cos C = \frac{3}{5}$．

（1）求 $\sin A$ 的值；

（2）若 $b=11$，求 $\triangle ABC$ 的面积．
\begin{solution}
\textbf{答案：}（1）$\dfrac{\sqrt{5}}{5}$；（2）$22$

\textbf{分析：}（1）先由平方关系求出 $\sin C$，再根据正弦定理即可解出；（2）根据余弦定理以及 $4a=\sqrt{5}c$ 可解出 $a$，即可由三角形面积公式求出面积．

\textbf{详解：}（1）由 $\cos C = \frac{3}{5}$，$0<C<\pi$，得 $\sin C = \frac{4}{5}$．由 $4a=\sqrt{5}c$ 及正弦定理得 $4\sin A = \sqrt{5}\sin C$，所以 $\sin A = \frac{\sqrt{5}}{4}\cdot\frac{4}{5} = \frac{\sqrt{5}}{5}$．
（2）由 $4a=\sqrt{5}c$ 设 $c = \frac{4}{\sqrt{5}}a$，由余弦定理 $\cos C = \frac{a^2+b^2-c^2}{2ab} = \frac{a^2+121-\frac{16}{5}a^2}{2\cdot a\cdot11} = \frac{3}{5}$，解得 $a=5$（负值舍去），所以 $\triangle ABC$ 的面积 $S = \frac{1}{2}ab\sin C = \frac{1}{2}\times5\times11\times\frac{4}{5}=22$．
\end{solution}
\end{question}

\begin{question}[points=15]
如图，已知 $ABCD$ 和 $CDEF$ 都是直角梯形，$AB\parallel DC, DC\parallel EF, AB=5, DC=3, EF=1, \angle BAD = \angle CDE = 60^\circ$，二面角 $F-DC-B$ 的平面角为 $60^\circ$．设 $M,N$ 分别为 $AE, BC$ 的中点．

（1）证明：$FN \perp AD$；

（2）求直线 $BM$ 与平面 $ADE$ 所成角的正弦值．
\begin{solution}
\textbf{答案：}（1）见解析；（2）$\dfrac{5\sqrt{7}}{14}$

\textbf{分析：}（1）过点 $E,D$ 分别做直线 $DC,AB$ 的垂线 $EG,DH$，由平面几何知识易得 $FC=BC$，再根据二面角的定义可知 $\angle BCF=60^\circ$，由此可证得 $FN\perp$ 平面 $ABCD$，即得 $FN\perp AD$；（2）建立空间直角坐标系，利用线面角的向量公式求解．

\textbf{详解：}（1）过 $E$ 作 $EG\perp DC$ 于 $G$，过 $D$ 作 $DH\perp AB$ 于 $H$，易得 $DG=AH=2$，$\angle EFC = \angle DCF = \angle DCB = \angle ABC = 90^\circ$，所以 $EG=DH=2\sqrt{3}$．由 $DC\perp CF, DC\perp CB$ 得 $DC\perp$ 平面 $BCF$，所以 $\angle BCF=60^\circ$，$\triangle BCF$ 是正三角形．$N$ 是 $BC$ 中点，故 $FN\perp BC$，又 $DC\perp$ 平面 $BCF$，$FN\subset$ 平面 $BCF$，故 $FN\perp CD$，所以 $FN\perp$ 平面 $ABCD$，则 $FN\perp AD$．
（2）以 $N$ 为原点，$NK$（平行 $AB$）、$NB$、$NF$ 分别为 $x,y,z$ 轴建立空间直角坐标系，则 $A(5,\sqrt{3},0), B(0,\sqrt{3},0), D(3,-\sqrt{3},0), E(1,0,3)$，$M\left(3,\frac{\sqrt{3}}{2},\frac{3}{2}\right)$，$\overrightarrow{BM} = \left(3,-\frac{\sqrt{3}}{2},\frac{3}{2}\right)$，$\overrightarrow{AD}=(-2,-2\sqrt{3},0)$，$\overrightarrow{DE}=(-2,\sqrt{3},3)$．设平面 $ADE$ 的法向量为 $\mathbf{n}=(x,y,z)$，由 $\mathbf{n}\cdot\overrightarrow{AD}=0$ 且 $\mathbf{n}\cdot\overrightarrow{DE}=0$ 得 $\begin{cases} -2x-2\sqrt{3}y=0 \\ -2x+\sqrt{3}y+3z=0 \end{cases}$，取 $\mathbf{n}=(\sqrt{3},-1,\sqrt{3})$．设 $BM$ 与平面 $ADE$ 所成角为 $\theta$，则 $\sin\theta = \frac{|\mathbf{n}\cdot\overrightarrow{BM}|}{|\mathbf{n}|\cdot|\overrightarrow{BM}|} = \frac{|3\sqrt{3}+\frac{\sqrt{3}}{2}+\frac{3\sqrt{3}}{2}|}{\sqrt{3+1+3}\cdot\sqrt{9+\frac{3}{4}+\frac{9}{4}}} = \frac{5\sqrt{7}}{14}$.
\end{solution}
\end{question}

\begin{question}[points=15]
已知等差数列 $\{a_n\}$ 的首项 $a_1 = -1$，公差 $d > 1$．记 $\{a_n\}$ 的前 $n$ 项和为 $S_n$ $(n \in \mathbb{N}^*)$．

（1）若 $S_4 - 2a_2 a_3 + 6 = 0$，求 $S_n$；

（2）若对于每个 $n \in \mathbb{N}^*$，存在实数 $c_n$，使 $a_n + c_n, a_{n+1} + 4c_n, a_{n+2} + 15c_n$ 成等比数列，求 $d$ 的取值范围．
\begin{solution}
\textbf{答案：}（1）$S_n = \dfrac{3n^2-5n}{2}$；（2）$1<d\le 2$

\textbf{分析：}（1）利用等差数列通项公式及前 $n$ 项和公式化简条件，求出 $d$，再求 $S_n$；（2）由等比数列定义列方程，结合一元二次方程有解的条件求 $d$ 的范围.

\textbf{详解：}（1）由 $S_4 - 2a_2a_3+6=0$，$a_1=-1$，得 $(-4+6d) - 2(-1+d)(-1+2d)+6=0$，化简得 $d^2-3d=0$，又 $d>1$，所以 $d=3$，$a_n=3n-4$，$S_n = \frac{n(a_1+a_n)}{2} = \frac{3n^2-5n}{2}$．
（2）由 $a_n+c_n, a_{n+1}+4c_n, a_{n+2}+15c_n$ 成等比数列得 $(a_{n+1}+4c_n)^2 = (a_n+c_n)(a_{n+2}+15c_n)$，代入 $a_n=-1+(n-1)d$，整理得 $c_n^2 + (14d-8nd+8)c_n + d^2 = 0$，关于 $c_n$ 的方程有实数解，判别式 $\Delta = (14d-8nd+8)^2 - 4d^2 \ge 0$ 对任意 $n\in\mathbb{N}^*$ 恒成立，即 $[(n-2)d-1][(2n-3)d-2] \ge 0$ 恒成立．当 $n=1$ 时，$(d+1)(d+2)\ge0$ 成立；当 $n=2$ 时，$(-1)(d-2)\ge0$ 得 $d\le2$；当 $n\ge3$ 时，$[(n-2)d-1][(2n-3)d-2] > (n-3)(2n-5) \ge 0$ 恒成立．又 $d>1$，所以 $1<d\le2$.
\end{solution}
\end{question}

\begin{question}[points=15]
如图，已知椭圆 $\dfrac{x^2}{12} + y^2 = 1$．设 $A,B$ 是椭圆上异于 $P(0,1)$ 的两点，且点 $Q\left(0,\frac{1}{2}\right)$ 在线段 $AB$ 上，直线 $PA, PB$ 分别交直线 $y = -\frac{1}{2}x + 3$ 于 $C, D$ 两点．

（1）求点 $P$ 到椭圆上点的距离的最大值；

（2）求 $|CD|$ 的最小值．
\begin{solution}
\textbf{答案：}（1）$\dfrac{12\sqrt{11}}{11}$；（2）$\dfrac{6\sqrt{5}}{5}$

\textbf{分析：}（1）设 $Q(2\sqrt{3}\cos\theta, \sin\theta)$ 是椭圆上任意一点，再根据两点间的距离公式求出 $|PQ|^2$，再根据二次函数的性质即可求出；（2）设直线 $AB: y=kx+\frac{1}{2}$ 与椭圆方程联立，可得 $x_1x_2, x_1+x_2$，再将直线 $y=-\frac{1}{2}x+3$ 与 $PA, PB$ 的方程分别联立，可解得点 $C,D$ 的坐标，再根据两点间的距离公式求出 $|CD|$，最后化简可得 $|CD| = \frac{3\sqrt{5}}{2}\cdot\frac{\sqrt{16k^2+1}}{3k+1}$，由柯西不等式即可求出最小值．

\textbf{详解：}（1）设椭圆上一点 $M(2\sqrt{3}\cos\theta,\sin\theta)$，则 $|PM|^2 = 12\cos^2\theta + (1-\sin\theta)^2 = 13 - 11\sin^2\theta - 2\sin\theta = -11\left(\sin\theta+\frac{1}{11}\right)^2 + \frac{144}{11} \le \frac{144}{11}$，当且仅当 $\sin\theta = -\frac{1}{11}$ 时取等，所以 $|PM|_{\max} = \frac{12\sqrt{11}}{11}$．
（2）设 $AB: y = kx + \frac{1}{2}$，代入椭圆方程得 $(k^2+\frac{1}{12})x^2 + kx - \frac{3}{4}=0$，设 $A(x_1,y_1), B(x_2,y_2)$，则 $x_1+x_2 = -\frac{k}{k^2+\frac{1}{12}}$，$x_1x_2 = -\frac{3}{4(k^2+\frac{1}{12})}$．$PA: y = \frac{y_1-1}{x_1}x+1$，与 $y=-\frac{1}{2}x+3$ 联立得 $x_C = \frac{4x_1}{x_1+2y_1-2} = \frac{4x_1}{(2k+1)x_1-1}$，同理 $x_D = \frac{4x_2}{(2k+1)x_2-1}$．则 $|CD| = \sqrt{1+\frac{1}{4}}|x_C-x_D| = \frac{\sqrt{5}}{2}\cdot\left|\frac{4x_1}{(2k+1)x_1-1} - \frac{4x_2}{(2k+1)x_2-1}\right| = 2\sqrt{5}\cdot\frac{|x_1-x_2|}{|(2k+1)^2 x_1x_2 - (2k+1)(x_1+x_2)+1|}$．代入韦达定理并化简得 $|CD| = \frac{3\sqrt{5}}{2}\cdot\frac{\sqrt{16k^2+1}}{3k+1}$．由柯西不等式 $(\sqrt{16k^2+1})^2(3^2+1^2) \ge (3\cdot\sqrt{16k^2+1}+1\cdot1)^2$，得 $\frac{\sqrt{16k^2+1}}{3k+1} \ge \frac{4}{5}$，当且仅当 $k=\frac{3}{16}$ 时取等，所以 $|CD|_{\min} = \frac{3\sqrt{5}}{2}\cdot\frac{4}{5} = \frac{6\sqrt{5}}{5}$.
\end{solution}
\end{question}

\begin{question}[points=15]
设函数 $f(x) = \dfrac{e}{2x} + \ln x$ $(x>0)$．

（1）求 $f(x)$ 的单调区间；

（2）已知 $a,b \in \mathbb{R}$，曲线 $y=f(x)$ 上不同的三点 $(x_1, f(x_1)), (x_2, f(x_2)), (x_3, f(x_3))$ 处的切线都经过点 $(a,b)$．证明：

（ⅰ）若 $a>e$，则 $0 < b - f(a) < \dfrac{1}{2}\left(\dfrac{a}{e} - 1\right)$；

（ⅱ）若 $0<a<e, x_1<x_2<x_3$，则 $\dfrac{2\sqrt{e-a}}{e} + \dfrac{1}{6e^2} < \dfrac{1}{x_1} + \dfrac{1}{x_3} < \dfrac{2\sqrt{e-a}}{a} - \dfrac{1}{6e}$．

（注：$e=2.71828\cdots$ 是自然对数的底数）
\begin{solution}
\textbf{答案：}见解析

\textbf{分析：}（1）求出导数，讨论其符号可得单调区间；（2）（ⅰ）由题设构造关于切点横坐标的方程，根据方程有3个不同的解可证明不等式成立；（ⅱ）通过比值代换等证明．

\textbf{详解：}（1）$f'(x) = -\frac{e}{2x^2} + \frac{1}{x} = \frac{2x-e}{2x^2}$，当 $0<x<\frac{e}{2}$ 时 $f'(x)<0$，当 $x>\frac{e}{2}$ 时 $f'(x)>0$，所以 $f(x)$ 的减区间为 $(0,\frac{e}{2})$，增区间为 $(\frac{e}{2},+\infty)$．
（2）（ⅰ）设切点 $(x,f(x))$，切线方程为 $y-f(x)=f'(x)(x-a)$，即 $g(x)=f(x)-b-f'(x)(x-a)=0$，化简得 $g(x)=1-\frac{a+e}{x}+\frac{ae}{2x^2}-\ln x+b$，$g'(x)=\frac{1}{x^3}(x-e)(x-a)$．当 $a>e$ 时，$g(x)$ 在 $(0,e)$、$(a,+\infty)$ 递减，在 $(e,a)$ 递增．由 $g(x)$ 有三个零点得 $g(e)<0$ 且 $g(a)>0$，代入得 $b<\frac{a}{2e}+1$ 且 $b>\frac{e}{2a}+\ln a=f(a)$，从而 $0<b-f(a)<\frac{1}{2}\left(\frac{a}{e}-1\right)$．
（ⅱ）类似（ⅰ）可得 $\frac{a}{2e}+1<b<\frac{e}{2a}+\ln a$．设 $t=\frac{e}{x}$，$m=\frac{a}{e}\in(0,1)$，则 $t_1,t_2,t_3$ 满足 $-\frac{m}{2}t^2+(m+1)t-\ln t-b=0$．由 $t_1t_3=1$ 及韦达定理可证．具体证明过程略．
\end{solution}
\end{question}

